\providecommand\SimplexChains{7}%
\providecommand\SimplexNodeColor{red}%
\providecommand\SimplexRadius{5}%
\providecommand\SimplexNodeRadius{3pt}%
% use this command to differentiate nodes from different simplexes
\providecommand\SimplexNodeName{chain}%
\providecommand\SimplexOffsetAngle{0}%
\providecommand\SimplexOffsetR{0}%
\providecommand\SimplexRotationMod{0}%
\providecommand\SimplexOutlineColor{black}%
\providecommand\SimplexLineScale{1}%
%
% ultra thick = 1.6pt
\pgfmathtruncatemacro\nodeLineThickness{1.6 * \SimplexLineScale}
\tikzstyle{refl} = [line width=\nodeLineThickness pt, draw=\SimplexOutlineColor, shorten >= 0pt, shorten <= 0pt]
\tikzstyle{splxNode} = [draw=\SimplexOutlineColor, line width=\nodeLineThickness pt, shape=circle]
%
\newcommand\TikzSimplex{
    \pgfmathsetmacro\rOffset{(360/\SimplexChains)+90+\SimplexRotationMod}
    \node (\SimplexNodeName-origin) at (\SimplexOffsetAngle: \SimplexOffsetR) {};
    % \node (\SimplexNodeName-origin-label) [yshift=\SimplexRadius*1.2cm] at (\SimplexNodeName-origin) {O \SimplexOffsetAngle: \SimplexOffsetR -- \SimplexNodeRadius};
    \foreach \i in {1,...,\SimplexChains} {
            \pgfmathsetmacro\r{\i*(360/\SimplexChains)+\rOffset}
            \xdef\sxRadius{\SimplexRadius}
            \ifnum 1=\SimplexChains
                \xdef\sxRadius{0}
            \fi
            \node (\SimplexNodeName//\i)
                [splxNode, fill=\SimplexNodeColor, draw=\SimplexOutlineColor, inner sep=\SimplexNodeRadius]
                at ([shift=({\r: \sxRadius*0.8cm})]\SimplexNodeName-origin)
                {};
        }
    \foreach \i in {1,...,\SimplexChains} {
            \foreach \j in {1,...,\SimplexChains} {
                    \ifnum \i>\j
                        \draw [refl, draw=\SimplexOutlineColor] (\SimplexNodeName//\i) -- (\SimplexNodeName//\j);
                        % \node (n-\i-\j) at (\i, \j) {\i, \j};
                    \else
                        \breakforeach
                    \fi
                }
        }
}%
